% =========================================================
% CONSIDERAÇÕES FINAIS
% FINALIDADE: responder ao problema de pesquisa e encerrar o trabalho.
% ESCREVA: a síntese das respostas, a retomada dos objetivos, as
%   contribuições, as limitações e as possibilidades de continuidade.
% Não introduza dados ou referências novas nesta seção.
% EXEMPLO NESTE ARQUIVO (comentado): remissão a uma seção anterior.
% =========================================================

\chapter{CONSIDERAÇÕES FINAIS}
\label{cap:consideracoes}

% Exemplo de remissão (apague o % do início da linha para usar):
% Conforme os objetivos definidos na Seção~\ref{sec:objetivos}, ...

As considerações finais retomam o problema de pesquisa e o objetivo geral e apresentam, de forma sintética, as respostas obtidas. Cada objetivo específico deve ser retomado, com a indicação dos resultados que permitiram alcançá-lo.

Em seguida, destacam-se as contribuições do estudo para a área de conhecimento, para a prática educativa e, quando for o caso, para as instituições ou as comunidades envolvidas. Quando a pesquisa resultar em produto educacional ou em procedimento aplicável, convém indicar também suas possibilidades de uso.

Por fim, reconhecem-se as limitações da pesquisa e indicam-se possibilidades de continuidade, como a ampliação do recorte investigado, o emprego de outros procedimentos metodológicos ou o aprofundamento de aspectos que não puderam ser tratados no âmbito deste trabalho.
